\documentclass[10pt,a4paper]{amsart}
\usepackage{amsmath,amssymb,amsthm,graphicx}
% Retained preamble for the Stacks Project chapter {CH} (shared project preamble).
% The chapter file pinned-original-source.tex only carries \input{{preamble}} on its first line,
% so no preamble text lives inside the chapter itself.  The Stacks Project uses one single shared
% preamble for every chapter: the snapshot stages it three times (sources/stacks-chow/preamble.tex,
% sources/stacks-etale-cohomology/preamble.tex and focused-seeds/AGP/sources/stacks-snapshot-20261011/
% preamble.tex) and those copies are identical in size.  The block below is copied mechanically,
% line for line, from sources/stacks-chow/preamble.tex lines 14-18 and 154-217: the Xy-pic setup and
% the native theorem declarations together with the mathematical macros.  Nothing was added, removed
% or redefined; the wrapper class and the amsmath/amssymb/amsthm/graphicx base packages are supplied
% by the assembler.
% For commutative diagrams we use Xy-pic
% Chapter: stacks-spaces-more-morphisms.tex -- More on Morphisms of Spaces.
\usepackage[all]{xy}

% We use 2cell for 2-commutative diagrams.
\xyoption{2cell}

% Theorem environments.
%
\theoremstyle{plain}
\newtheorem{theorem}[subsection]{Theorem}
\newtheorem{proposition}[subsection]{Proposition}
\newtheorem{lemma}[subsection]{Lemma}

\theoremstyle{definition}
\newtheorem{definition}[subsection]{Definition}
\newtheorem{example}[subsection]{Example}
\newtheorem{exercise}[subsection]{Exercise}
\newtheorem{situation}[subsection]{Situation}

\theoremstyle{remark}
\newtheorem{remark}[subsection]{Remark}
\newtheorem{remarks}[subsection]{Remarks}

\numberwithin{equation}{subsection}

% Macros
%
\def\lim{\mathop{\mathrm{lim}}\nolimits}
\def\colim{\mathop{\mathrm{colim}}\nolimits}
\def\Spec{\mathop{\mathrm{Spec}}}
\def\Hom{\mathop{\mathrm{Hom}}\nolimits}
\def\Ext{\mathop{\mathrm{Ext}}\nolimits}
\def\SheafHom{\mathop{\mathcal{H}\!\mathit{om}}\nolimits}
\def\SheafExt{\mathop{\mathcal{E}\!\mathit{xt}}\nolimits}
\def\Sch{\mathit{Sch}}
\def\Mor{\mathop{\mathrm{Mor}}\nolimits}
\def\Ob{\mathop{\mathrm{Ob}}\nolimits}
\def\Sh{\mathop{\mathit{Sh}}\nolimits}
\def\NL{\mathop{N\!L}\nolimits}
\def\CH{\mathop{\mathrm{CH}}\nolimits}
\def\proetale{{pro\text{-}\acute{e}tale}}
\def\etale{{\acute{e}tale}}
\def\QCoh{\mathit{QCoh}}
\def\Ker{\mathop{\mathrm{Ker}}}
\def\Im{\mathop{\mathrm{Im}}}
\def\Coker{\mathop{\mathrm{Coker}}}
\def\Coim{\mathop{\mathrm{Coim}}}

% Boxtimes
%
\DeclareMathSymbol{\boxtimes}{\mathbin}{AMSa}{"02}

%
% Macros for moduli stacks/spaces
%
\def\QCohstack{\mathcal{QC}\!\mathit{oh}}
\def\Cohstack{\mathcal{C}\!\mathit{oh}}
\def\Spacesstack{\mathcal{S}\!\mathit{paces}}
\def\Quotfunctor{\mathrm{Quot}}
\def\Hilbfunctor{\mathrm{Hilb}}
\def\Curvesstack{\mathcal{C}\!\mathit{urves}}
\def\Polarizedstack{\mathcal{P}\!\mathit{olarized}}
\def\Complexesstack{\mathcal{C}\!\mathit{omplexes}}
% \Pic is the operator that assigns to X its picard group, usage \Pic(X)
% \Picardstack_{X/B} denotes the Picard stack of X over B
% \Picardfunctor_{X/B} denotes the Picard functor of X over B
\def\Pic{\mathop{\mathrm{Pic}}\nolimits}
\def\Picardstack{\mathcal{P}\!\mathit{ic}}
\def\Picardfunctor{\mathrm{Pic}}
\def\Deformationcategory{\mathcal{D}\!\mathit{ef}}
\newenvironment{reference}{}{}
\newenvironment{slogan}{}{}
\newenvironment{history}{}{}
\usepackage{cleveref}
\crefname{originallemma}{original lemma}{original lemmas}
\Crefname{originallemma}{Original lemma}{Original lemmas}
\makeatletter
\newlabel{algebra-definition-small-extension}{{algebra-definition-small-extension}{source}{}{sourcepointer-1}{}}
\newlabel{algebra-definition-small-extension@cref}{{[originallemma][0][]algebra-definition-small-extension}{[1][0][]source}{}{}{}}
\makeatother
\crefname{originallemma}{original lemma}{original lemmas}
\Crefname{originallemma}{Original lemma}{Original lemmas}
\makeatletter
\newlabel{lemma-etale-formally-etale}{{lemma-etale-formally-etale}{source}{}{sourcepointer-2}{}}
\newlabel{lemma-etale-formally-etale@cref}{{[originallemma][0][]lemma-etale-formally-etale}{[1][0][]source}{}{}{}}
\makeatother
\crefname{originallemma}{original lemma}{original lemmas}
\Crefname{originallemma}{Original lemma}{Original lemmas}
\makeatletter
\newlabel{lemma-smooth-formally-smooth}{{lemma-smooth-formally-smooth}{source}{}{sourcepointer-3}{}}
\newlabel{lemma-smooth-formally-smooth@cref}{{[originallemma][0][]lemma-smooth-formally-smooth}{[1][0][]source}{}{}{}}
\makeatother
\crefname{originallemma}{original lemma}{original lemmas}
\Crefname{originallemma}{Original lemma}{Original lemmas}
\makeatletter
\newlabel{more-morphisms-lemma-lifting-along-artinian-at-point}{{more-morphisms-lemma-lifting-along-artinian-at-point}{source}{}{sourcepointer-4}{}}
\newlabel{more-morphisms-lemma-lifting-along-artinian-at-point@cref}{{[originallemma][0][]more-morphisms-lemma-lifting-along-artinian-at-point}{[1][0][]source}{}{}{}}
\makeatother
\title{Smoothness at a point of a morphism of algebraic spaces is equivalent to lifting along small extensions there}
\author{The Stacks Project Authors}
\date{stacks-project:spaces-more-morphisms.tex:local-snapshot-20261011; GNU FDL 1.2 or later}
\begin{document}
\maketitle
\noindent\textit{Editorial note.} Statement span is lines 5095-5117 of sources/stacks-spaces-more-morphisms/pinned-original-source.tex; proof span is lines 5119-5152, the complete author proof (32 body lines).\par\medskip
\noindent\textit{Editorial note.} Navigation stubs: 4 cited-but-not-included labels are supplied as non-load-bearing pointers with their pinned tags where the snapshot tag index has them.\par\medskip
\noindent\textit{Editorial note.} extra\_packages is empty on purpose: the retained preamble already loads xy with options and a second declaration clashes in the assembler wrapper.\par\medskip
\noindent\textit{Editorial note.} Text is retained verbatim, including the author's own typos and placeholders.\par\medskip
\noindent\textit{Allowed standard prerequisites (editorial).}\begin{itemize}
\item Explicit prerequisite boundary. Smoothness at a point, small extensions of local rings, and the schemes-side Artinian lifting criterion; the reduction is retained.
\end{itemize}\medskip
\ifdefined\hypertarget\hypertarget{sourcepointer-1}{}\fi
\noindent\textit{Source reference.} Original lemma algebra-definition-small-extension: sources/stacks-spaces-more-morphisms/pinned-original-source.tex, source lines 5152--5152. Cited by the retained statement or proof but not part of the collected spans; supplied only as a navigation pointer. The argument does not depend on its content.\par\medskip
\ifdefined\hypertarget\hypertarget{sourcepointer-2}{}\fi
\noindent\textit{Source reference.} Original lemma lemma-etale-formally-etale: sources/stacks-spaces-more-morphisms/pinned-original-source.tex, source lines 3641--3649. Cited by the retained statement or proof but not part of the collected spans; supplied only as a navigation pointer. The argument does not depend on its content.\par\medskip
\ifdefined\hypertarget\hypertarget{sourcepointer-3}{}\fi
\noindent\textit{Source reference.} Original lemma lemma-smooth-formally-smooth: sources/stacks-spaces-more-morphisms/pinned-original-source.tex, source lines 4587--4595. Cited by the retained statement or proof but not part of the collected spans; supplied only as a navigation pointer. The argument does not depend on its content.\par\medskip
\ifdefined\hypertarget\hypertarget{sourcepointer-4}{}\fi
\noindent\textit{Source reference.} Original lemma more-morphisms-lemma-lifting-along-artinian-at-point: sources/stacks-spaces-more-morphisms/pinned-original-source.tex, source lines 5152--5152. Cited by the retained statement or proof but not part of the collected spans; supplied only as a navigation pointer. The argument does not depend on its content.\par\medskip
\begin{lemma}
\label{lemma-lifting-along-artinian-at-point}
Let $S$ be a scheme. Let $f : X \to Y$ be a morphism of algebraic spaces
over $S$. Let $x \in |X|$.
Assume that $Y$ is locally Noetherian and $f$ locally of finite type.
The following are equivalent:
\begin{enumerate}
\item $f$ is smooth at $x$,
\item for every solid commutative diagram
$$
\xymatrix{
X \ar[d]_f & \Spec(B) \ar[d]^i \ar[l]^-\alpha \\
Y & \Spec(B') \ar[l]_-{\beta} \ar@{-->}[lu]
}
$$
where $B' \to B$ is a surjection of local rings with
$\Ker(B' \to B)$ of square zero, and $\alpha$ mapping the
closed point of $\Spec(B)$ to $x$ there exists
a dotted arrow making the diagram commute, and
\item same as in (2) but with $B' \to B$ ranging over small
extensions (see Algebra, Definition \ref{algebra-definition-small-extension}).
\end{enumerate}
\end{lemma}

\begin{proof}
Condition (1) means there is an open subspace $X' \subset X$
such that $X' \to Y$ is smooth. Hence (1) implies conditions (2) and (3) by
Lemma \ref{lemma-smooth-formally-smooth}. Condition (2) implies
condition (3) trivially. Assume (3). Choose a commutative diagram
$$
\xymatrix{
X \ar[d] & U \ar[l] \ar[d] \\
Y & V \ar[l]
}
$$
with $U$ and $V$ affine, horizontal arrows \'etale and
such that there is a point $u \in U$ mapping to $x$. Next, consider
a diagram
$$
\xymatrix{
X \ar[d] & U \ar[l] \ar[d] & \Spec(B) \ar[d]^i \ar[l]^-\alpha  \\
Y & V \ar[l] & \Spec(B') \ar[l]_-{\beta}
}
$$
as in (3) but for $u \in U \to V$. Let $\gamma : \Spec(B') \to X$
be the arrow we get from our assumption that (3) holds for $X$.
Because $U \to X$ is \'etale and
hence formally \'etale (Lemma \ref{lemma-etale-formally-etale})
the morphism $\gamma$
has a unique lift to $U$ compatible with $\alpha$. Then because
$V \to Y$ is \'etale hence formally \'etale this lift is compatible
with $\beta$. Hence (3) holds for $u \in U \to V$ and we conclude
that $U \to V$ is smooth at $u$ by
More on Morphisms, Lemma
\ref{more-morphisms-lemma-lifting-along-artinian-at-point}.
This proves that $X \to Y$ is smooth at $x$, thereby
finishing the proof.
\end{proof}

\end{document}
